\section{Modos de falla y líneas de defensa de ingeniería}

\subsection{Falla típica: imposibilidad de recuperarse tras una acción errónea}
El ponente puso, durante la sesión de preguntas y respuestas, un ejemplo muy representativo: cuando el modelo ejecuta una acción errónea y carece de búsqueda y retroceso, resulta muy difícil autocorregirse después, y termina por desviarse por completo del objetivo de la tarea.

\begin{figure}[H]
\centering
\includegraphics[width=0.82\textwidth]{frame-09-failure-mode.jpg}
\caption{Video en aproximadamente 01:17:20: el ponente discute la dificultad de recuperarse tras una acción errónea y enfatiza la necesidad de mecanismos de búsqueda y retroceso}
\end{figure}

\begin{importantbox}{El ciclo mínimo de recuperación ante fallas}
\begin{enumerate}
    \item Detección: identificar rápidamente que la trayectoria actual se desvió;
    \item Reversión: volver al estado confiable más reciente;
    \item Replanificación: actualizar los subobjetivos y las restricciones de acción;
    \item Revisión: determinar mediante el verifier si se volvió a la trayectoria correcta.
\end{enumerate}
\end{importantbox}

\subsection{Clasificación de los modos de falla}
\begin{table}[H]
\centering
\small
\begin{tabular}{p{2.7cm}p{3.8cm}p{4.2cm}}
\toprule
\textbf{Tipo de falla} & \textbf{Síntomas típicos} & \textbf{Medidas de corrección de ingeniería} \\
\midrule
Error de percepción & Lectura equivocada de botones, campos de formulario o contenido de imágenes & Verificación multimodal, umbral de confianza de elementos \\
Error de planificación & Orden equivocado de subobjetivos u omisión de restricciones & Plantillas de plan, verificador de subobjetivos \\
Error de ejecución & Clics, entradas o detenciones equivocadas de manera continua & Registro de repetición, instantáneas de estado y reversión \\
Error de evaluación & El Judge da un veredicto equivocado de éxito o fracaso & Cruce entre varios Judges, calibración mediante revisión manual por muestreo \\
\bottomrule
\end{tabular}
\caption{Modos de falla comunes del Web Agent y estrategias de corrección}
\end{table}

\begin{warningbox}{No trates las ``fallas'' como muestras de ruido}
Las trayectorias de falla son una de las fuentes de datos más valiosas. El sistema debe archivar de manera sistemática los tipos de falla y realimentarlos al flujo de entrenamiento y evaluación, formando un ciclo de mejora continua.
\end{warningbox}

\subsection{La colaboración humana sigue siendo necesaria}
Para tareas de alto riesgo, el agente debe tener la capacidad de ``solicitar aclaración'' y de ``escalar a una persona''. La persona no debe aparecer solo como el último respaldo, sino intervenir de manera oportuna cuando aumenta la incertidumbre.

\begin{knowledgebox}{Ejemplos de condiciones que activan la colaboración humano-máquina}
\begin{itemize}
    \item Dos discrepancias consecutivas de alta confianza (el planner y el verifier no coinciden);
    \item Un campo clave no se puede determinar (relacionado con pagos, identidad o privacidad);
    \item Se rebasa el presupuesto pero la tarea no converge.
\end{itemize}
\end{knowledgebox}

\subsection{Resumen de la sección}
Los modos de falla del agente pueden gestionarse de manera sistemática. Un sistema verdaderamente confiable no es aquel que ``nunca se equivoca'', sino uno cuyos errores son ``detectables, recuperables y aprovechables para el aprendizaje''.

